% !TEX program = pdflatex
\documentclass[12pt, a4paper]{article}

\usepackage[a4paper, margin=2.5cm]{geometry}
\usepackage[utf8]{inputenc}
\usepackage[T1]{fontenc}
\usepackage{mathptmx} % Font Times New Roman ukuran 12pt

\usepackage{enumitem}
\setlist[itemize]{label=-}
\usepackage{amsmath}
\usepackage{booktabs}
\usepackage[table]{xcolor}
\usepackage{listings}
\usepackage{hyperref}
\usepackage{graphicx}
\graphicspath{ {images/} } % Opsional: Menentukan folder tempat Anda menyimpan semua foto


% Mengubah label bahasa Inggris ke Indonesia tanpa perlu package babel
\AtBeginDocument{%
    \renewcommand{\tablename}{Tabel}%
    \renewcommand{\figurename}{Gambar}%
    \renewcommand{\lstlistingname}{Kode}%
}

\definecolor{codebg}{rgb}{0.96,0.96,0.96}
\definecolor{codeframe}{rgb}{0.85,0.85,0.85}

\lstset{
    backgroundcolor=\color{codebg},
    frame=single,
    rulecolor=\color{codeframe},
    basicstyle=\ttfamily\small,
    breaklines=true,
    captionpos=b,
    numbers=left,
    numberstyle=\tiny\color{gray},
    showstringspaces=false,
    tabsize=4    
}

\hypersetup{
    colorlinks=true,
    linkcolor=blue,
    urlcolor=blue,
    citecolor=blue
}

\begin{document}

% 1. HALAMAN JUDUL
\begin{titlepage}
    \centering
    \vspace*{3cm}
    {\large \textbf{LAPORAN UJIAN TENGAH SEMESTER}\\}   
    \vspace{0.5cm} 
    {\large \textbf{Sistem Tertanam}\\}
    \vspace{0.5cm}
    {\large "Monitoring Ruangan Sederhana Menggunakan Raspberry Pi 5 dengan Sensor DHT11, PIR, dan Buzzer"\\}
    \vspace{1cm}
    {\large \textbf{Dosen Pengampu:}\\}
    {\large \textbf{Drs. Hari Suparwito, S.J., M.App.IT}}\\
    \vspace{1cm}

    \includegraphics[width=7cm]{usd.png}\\[0.5cm]

    \vfill
    \text{Disusun Oleh:}\\
    \vspace{0.3cm}
    {\large \textbf{Filipus Ferry Pradana / 235314039}}\\
    {\large \textbf{I Dewa Gede Khrisnayana Sathya Gautama / 235314114}}\\

    \vspace{1cm}
    {\large \textbf{PROGRAM STUDI INFORMATIKA}}\\
    \vspace{0.3cm}
    {\large \textbf{FAKULTAS SAINS DAN TEKNOLOGI}}\\
    \vspace{0.3cm}
    {\large \textbf{UNIVERSITAS SANATA DHARMA}}\\
    \vspace{0.3cm}
    {\large \textbf{2026}}
    \vspace{0.3cm}
\end{titlepage}



% 2. PENDAHULUAN
\section{Pendahuluan}
\subsection{Latar Belakang}
Kondisi lingkungan dan keamanan fisik di dalam ruangan khusus, seperti ruang laboratorium atau ruang server, memegang peranan 
krusial dalam menjaga keandalan serta usia pakai perangkat elektronik di dalamnya. Suhu ruangan yang terlalu tinggi atau tingkat 
kelembaban udara yang ekstrem dapat mempercepat degradasi komponen perangkat keras, memicu korsleting akibat kondensasi, dan 
menurunkan performa sistem. Di samping faktor termal dan kelembaban, pemantauan aktivitas atau pergerakan manusia di dalam ruangan 
juga diperlukan sebagai lapisan keamanan fisik untuk mendeteksi keberadaan personel maupun potensi akses tidak sah. Oleh karena itu,
 diperlukan sebuah sistem berbasis \textit{embedded system} yang mampu memantau parameter suhu, kelembaban udara, serta keberadaan 
 gerakan secara \textit{real-time} menggunakan sensor DHT11 dan sensor PIR (\textit{Passive Infrared}), sekaligus memberikan 
 peringatan bertingkat melalui indikator visual (LED) dan audio (\textit{buzzer}).

\subsection{Rumusan Masalah}
Berdasarkan latar belakang di atas, rumusan masalah pada proyek ini adalah:
\begin{enumerate}
    \item Bagaimana merancang dan mengintegrasikan sensor suhu dan kelembaban (DHT11), sensor deteksi gerakan (PIR), tombol 
    interupsi (\textit{push button}), serta indikator output dengan Raspberry Pi?
    \item Bagaimana mengimplementasikan logika \textit{Finite State Machine} (FSM) untuk mengelola masa kalibrasi 
    (\textit{warm-up}) sensor PIR serta menentukan prioritas peringatan berdasarkan ambang batas suhu, kelembaban, dan deteksi 
    gerakan?
    \item Bagaimana kinerja dan keandalan sistem dalam merespons perubahan kondisi lingkungan, deteksi aktivitas gerakan, serta 
    penekanan tombol peredam alarm (\textit{mute})?
\end{enumerate}

\subsection{Tujuan Proyek}
Tujuan dari pelaksanaan proyek ini adalah:
\begin{enumerate}
    \item Membangun purwarupa sistem monitoring ruangan sederhana berbasis Raspberry Pi yang mengintegrasikan pemantauan kondisi 
    termal/kelembaban dan deteksi gerakan.
    \item Mengimplementasikan pembacaan data sensor DHT11 dan sensor PIR secara berkala menggunakan pemrograman Python serta 
    mengatur transisi indikator output (LED \textit{Traffic Light} dan \textit{Buzzer}) secara akurat.
    \item Menguji responsivitas sistem terhadap berbagai skenario uji, meliputi masa \textit{warm-up} sensor PIR, kondisi normal, 
    peringatan waspada (kelembaban meningkat atau terdeteksi gerakan), kondisi kritis, hingga fitur peredaman suara alarm 
    (\textit{mute}).
\end{enumerate}

% 3. TINJAUAN PUSTAKA
\section{Tinjauan Pustaka}
\subsection{Raspberry Pi}
Raspberry Pi adalah komputer papan tunggal (\textit{Single Board Computer} / SBC) seukuran kartu kredit yang dilengkapi dengan
 prosesor berbasis arsitektur ARM, memori RAM, port komunikasi, serta deretan pin GPIO. Raspberry Pi umumnya menjalankan sistem 
 operasi berbasis Linux (seperti Raspberry Pi OS) dan mendukung pemrograman tingkat tinggi menggunakan bahasa Python untuk 
 mengendalikan perangkat keras eksternal secara efisien.

\subsection{Sensor Suhu dan Kelembaban (DHT11)}
Sensor DHT11 merupakan modul sensor digital yang berfungsi mengukur suhu dan kelembaban relatif udara di sekitarnya. Sensor ini 
menggunakan elemen sensor kapasitif untuk mengukur kelembaban dan termistor (\textit{Negative Temperature Coefficient}) untuk 
mengukur suhu, kemudian mengirimkan sinyal digital terkalibrasi melalui antarmuka komunikasi satu jalur (\textit{single-wire 
serial interface}) ke pin GPIO Raspberry Pi.

\subsection{Sensor Gerak PIR (\textit{Passive Infrared})}
Sensor PIR (\textit{Passive Infrared}) adalah sensor elektronik yang berfungsi mendeteksi pergerakan objek berdasarkan pancaran 
radiasi sinar inframerah pasif (panas tubuh) dari manusia atau benda hidup di sekitarnya. Istilah \textit{passive} menunjukkan 
bahwa sensor ini tidak memancarkan sinyal inframerah sendiri, melainkan hanya menerima radiasi termal melalui lensa Fresnel menuju 
elemen sensor piroelektrik di dalamnya. Ketika terjadi perubahan pola radiasi inframerah akibat gerakan, sensor PIR akan mengubah
 logika pin keluaran digitalnya menjadi \textit{HIGH} (3.3V). Saat pertama kali dialiri daya, elemen piroelektrik pada sensor PIR 
 memerlukan waktu stabilisasi atau kalibrasi termal (\textit{warm-up time}) sekitar 30 hingga 60 detik terhadap suhu latar ruangan
  agar dapat mendeteksi gerakan secara akurat tanpa memicu pembacaan palsu (\textit{false trigger}).

\subsection{Konsep Embedded System}
Sistem tertanam (\textit{embedded system}) adalah sistem komputasi khusus yang menggabungkan perangkat keras (\textit{hardware})
 dan perangkat lunak (\textit{software}) yang dirancang untuk menjalankan fungsi pengawasan atau pengendalian spesifik. 
 Sistem ini bekerja secara berkesinambungan dengan siklus utama menerima masukan dari berbagai sensor (\textit{input}), 
 memproses data sesuai algoritma atau \textit{state machine} yang ditanamkan (\textit{process}), dan menghasilkan aksi pada 
 aktuator atau indikator (\textit{output}).

\subsection{Konsep Input-Output GPIO}
GPIO (\textit{General Purpose Input/Output}) adalah pin antarmuka fisik pada Raspberry Pi yang dapat dikonfigurasi secara 
fleksibel melalui perangkat lunak sebagai jalur masukan (\textit{input}) maupun keluaran (\textit{output}) sinyal digital. 
Sebagai \textit{input}, pin GPIO membaca level logika tegangan tinggi (\textit{HIGH} / 3.3V) atau rendah (\textit{LOW} / 0V), 
seperti sinyal deteksi dari sensor PIR, untai data serial dari DHT11, maupun penekanan \textit{push button}. Sebagai
 \textit{output}, pin GPIO mengirimkan sinyal tegangan untuk mengaktifkan komponen indikator seperti modul LED dan \textit{buzzer}.

% 4. PERANCANGAN SISTEM
\section{Perancangan Sistem}
\subsection{Arsitektur Sistem}
Secara garis besar, sistem monitoring ruangan sederhana ini terbagi menjadi tiga bagian utama: bagian masukan (\textit{Input}),
 bagian pemrosesan (\textit{Process}), dan bagian keluaran (\textit{Output}). Diagram blok arsitektur sistem ditunjukkan pada 
 skema di bawah ini:

\vspace{0.3cm}
\noindent\framebox[\textwidth]{%
\parbox{0.95\textwidth}{%
\centering
\vspace{0.3cm}
\textbf{DIAGRAM BLOK ARSITEKTUR SISTEM}\\[0.3cm]
\textbf{[INPUT]} Sensor Suhu \& Kelembaban (DHT11) + Sensor Gerak (PIR) + Push Button\\
$\downarrow$\\
\textbf{[PROSES]} Raspberry Pi (Akuisisi Data \& Evaluasi Ambang Batas)\\
$\downarrow$\\
\textbf{[OUTPUT]} Indikator LED / Buzzer + Tampilan Terminal / Layar
\vspace{0.3cm}
}}
\vspace{0.3cm}

\subsection{Desain Software dan Hardware}
\subsubsection{Skema Rangkaian}
Berikut adalah representasi skema perancangan rangkaian perangkat keras yang menghubungkan sensor, \textit{push button}, dan 
indikator ke Raspberry Pi:

\vspace{0.2cm}
\noindent\framebox[\textwidth]{%
\parbox{0.95\textwidth}{%
\centering
\vspace{0.8cm}
\textit{[Tempatkan penjelasan detail rangkaian atau pengganti gambar skema rangkaian di sini]}\\[0.2cm]
\small (Catatan: Sensor DHT dihubungkan ke VCC 3.3V, GND, dan Pin Data GPIO; Push Button dihubungkan ke Pin GPIO dengan
 konfigurasi Pull-Up/Pull-Down; LED dihubungkan seri dengan resistor 220 $\Omega$ menuju Pin GPIO dan GND).
\vspace{0.8cm}
}}
\vspace{0.2cm}

\subsubsection{Penjelasan Pin GPIO}
Pemetaan alokasi pin GPIO pada Raspberry Pi (menggunakan penomoran BCM) yang digunakan dalam sistem monitoring ruang 
lab/server ini disajikan pada Tabel \ref{tab:pin_gpio}.

\begin{table}[h!]
\centering
\caption{Pemetaan Pin GPIO Raspberry Pi}
\vspace{0.2cm}
\begin{tabular}{@{}llll@{}}
\toprule
\textbf{Komponen} & \textbf{Pin BCM (GPIO)} & \textbf{Library / Objek} & \textbf{Fungsi / Mode} \\ \midrule
Sensor DHT11 (Data) & GPIO 4 (\texttt{board.D4}) & \texttt{adafruit\_dht.DHT11} & Input Data Suhu \& Lembab \\
Sensor PIR & GPIO 23 & \texttt{MotionSensor} & Input Deteksi Gerakan \\
Push Button (Mute) & GPIO 5 & \texttt{Button} (Pull-Up) & Input Interupsi Mute Alarm \\
LED Merah & GPIO 17 & \texttt{TrafficLights (red)} & Output Indikator Kritis \\
LED Kuning & GPIO 27 & \texttt{TrafficLights (yellow)} & Output Indikator Waspada \\
LED Hijau & GPIO 22 & \texttt{TrafficLights (green)} & Output Indikator Normal/Init \\
Active Buzzer & GPIO 25 & \texttt{Buzzer} & Output Alarm Audio \\ \bottomrule
\end{tabular}
\label{tab:pin_gpio}
\end{table}

\subsubsection{State Machine}
Sistem dirancang menggunakan pendekatan \textit{Finite State Machine} (FSM) dengan lima kondisi (\textit{state}) 
utama yang diatur berdasarkan skala prioritas kedaruratan seperti pada Tabel \ref{tab:state_machine}.

\begin{table}[h!]
\centering
\caption{Transisi State Machine Sistem Monitoring Ruangan}
\vspace{0.2cm}
\resizebox{\textwidth}{!}{%
\begin{tabular}{@{}llll@{}}
\toprule
\textbf{State} & \textbf{Kondisi Pemicu (Trigger)} & \textbf{Prioritas} & \textbf{Aksi Output (LED \& Buzzer)} \\ \midrule
\texttt{INITIALIZING} & Waktu \textit{warm-up} PIR $<$ 60 detik \& Suhu/RH aman & Awal & LED Hijau ON, Buzzer OFF \\
\texttt{NORMAL} & Suhu $<$ 31$^\circ$C, RH $<$ 75\%, Tidak ada gerakan & Normal & LED Hijau ON, Buzzer OFF \\
\texttt{WARNING} & RH $\ge$ 75\% ($<$ 80\%) ATAU terdeteksi gerakan PIR & Prioritas 2 & LED Kuning ON, Buzzer Beep Lambat (0.2s ON / 0.8s OFF) \\
\texttt{CRITICAL} & Suhu $\ge$ 31$^\circ$C ATAU RH $\ge$ 80\% & Prioritas 1 & LED Merah Kedip Cepat (0.2s), Buzzer Beep Cepat (0.1s) \\
\texttt{MUTED} & Tombol ditekan saat \texttt{WARNING} / \texttt{CRITICAL} & Override Audio & LED Kuning \& Merah ON, Buzzer OFF \\ \bottomrule
\end{tabular}%
}
\label{tab:state_machine}
\end{table}

\subsubsection{Flowchart Sistem}
Alur kerja perangkat lunak pada sistem monitoring ruangan ini terdiri dari tahapan berurutan berikut:
\begin{enumerate}
    \item \textbf{Mulai (Start):} Sistem dinyalakan dan program dijalankan.
    \item \textbf{Inisialisasi:} Mengatur mode penomoran pin GPIO (BCM), mengatur arah pin (input/output), dan menginisialisasi 
    objek sensor DHT.
    \item \textbf{Baca Input:} Membaca status \textit{push button} serta mengakuisisi nilai suhu ($^\circ$C) dan kelembaban 
    (\%) dari sensor DHT.
    \item \textbf{Evaluasi Kondisi:}
    \begin{itemize}
        \item Jika tombol ditekan $\rightarrow$ Reset alarm / tampilkan status manual di layar.
        \item Jika Suhu $>$ [Batas Suhu, misal 30$^\circ$C] ATAU Kelembaban $>$ [Batas Kelembaban, misal 75\%] $\rightarrow$
         Masuk ke kondisi \textbf{Peringatan (Warning)}: Nyalakan LED Merah/Buzzer, matikan LED Hijau.
        \item Jika Suhu dan Kelembaban berada di bawah batas $\rightarrow$ Masuk ke kondisi \textbf{Normal}: Nyalakan LED Hijau, 
        matikan LED Merah/Buzzer.
    \end{itemize}
    \item \textbf{Tampilkan Data:} Mencetak nilai suhu, kelembaban, dan status ruangan ke terminal.
    \item \textbf{Delay \& Loop:} Tunggu selama 2 detik, lalu kembali ke langkah \textbf{Baca Input} hingga program dihentikan 
    pengguna.
\end{enumerate}

\subsubsection{State Machine}
Sistem dirancang menggunakan pendekatan \textit{Finite State Machine} (FSM) dengan empat kondisi (\textit{state}) utama 
seperti pada Tabel \ref{tab:state_machine}.

\begin{table}[h!]
\centering
\caption{Transisi State Machine Sistem Monitoring Ruangan}
\vspace{0.2cm}
\begin{tabular}{@{}llll@{}}
\toprule
\textbf{State Saat Ini} & \textbf{Kondisi Pemicu (Trigger)} & \textbf{State Berikutnya} & \textbf{Aksi Output} \\ \midrule
\texttt{IDLE / INIT} & Inisialisasi selesai & \texttt{NORMAL} & Semua LED mati \\
\texttt{NORMAL} & Suhu $\le$ 30$^\circ$C \& RH $\le$ 75\% & \texttt{NORMAL} & LED Hijau ON, Alarm OFF \\
\texttt{NORMAL} & Suhu $>$ 30$^\circ$C atau RH $>$ 75\% & \texttt{WARNING} & LED Hijau OFF, Alarm ON \\
\texttt{WARNING} & Push Button ditekan & \texttt{MUTE / ACK} & LED Merah berkedip, Buzzer OFF \\
\texttt{WARNING / MUTE} & Suhu $\le$ 30$^\circ$C \& RH $\le$ 75\% & \texttt{NORMAL} & LED Hijau ON, Alarm OFF \\ \bottomrule
\end{tabular}
\label{tab:state_machine}
\end{table}

% 5. IMPLEMENTASI
% 5. IMPLEMENTASI
\section{Implementasi}
\subsection{Tools \& Library}
Perangkat lunak dan pustaka (\textit{library}) yang digunakan dalam tahap implementasi proyek ini meliputi:
\begin{itemize}
    \item \textbf{Bahasa Pemrograman:} Python 3 pada lingkungan sistem operasi berbasis Linux (Raspberry Pi OS).
    \item \textbf{Library \texttt{gpiozero}:} Menggunakan kelas abstraksi perangkat keras tingkat tinggi meliputi 
    \texttt{TrafficLights} (mengelola tiga indikator LED sekaligus), \texttt{Buzzer}, \texttt{Button} (dengan fitur 
    \textit{internal pull-up} dan \textit{software debouncing}), serta \texttt{MotionSensor} untuk sensor PIR.
    \item \textbf{Library \texttt{adafruit\_dht} \& \texttt{board}:} Pustaka CircuitPython untuk mengakuisisi data 
    digital dari sensor DHT11 melalui pin \texttt{board.D4} dengan konfigurasi \texttt{use\_pulseio=False} agar stabil di 
    Raspberry Pi.
    \item \textbf{Library \texttt{time}:} Menggunakan \texttt{time.monotonic()} untuk penghitungan waktu kalibrasi 
    (\textit{warm-up}) sensor PIR yang bebas dari gangguan perubahan jam sistem, serta \texttt{time.sleep()} untuk interval 
    \textit{polling} siklus utama.
\end{itemize}

\subsection{Struktur Program}
Program dibagi ke dalam beberapa blok fungsional yang modular, yaitu:
\begin{enumerate}
    \item \textbf{Inisialisasi Hardware \& Konfigurasi Sistem:} Mendefinisikan objek komponen beserta nomor pin GPIO (BCM) serta 
    menetapkan konstanta ambang batas (\texttt{BATAS\_SUHU\_TINGGI = 31.0}$^\circ$C, \texttt{BATAS\_LEMBAB\_WASPADA = 75.0}\%, 
    \texttt{BATAS\_LEMBAB\_KRITIS = 80.0}\%, dan \texttt{WAKTU\_WARMUP\_PIR = 60} detik).
    \item \textbf{Handler Interupsi Tombol (\texttt{tekan\_tombol\_handler}):} Fungsi \textit{callback} asinkron yang terhubung ke
     \texttt{tombol\_mute.when\_pressed}. Fungsi ini mengubah status \texttt{is\_muted} secara \textit{toggle} hanya ketika sistem 
     sedang berada dalam kondisi \texttt{WARNING}, \texttt{CRITICAL}, atau \texttt{MUTED}.
    \item \textbf{Manajemen Kalibrasi PIR (\texttt{periksa\_warmup\_pir}):} Menghitung mundur waktu stabilisasi elemen piroelektrik
     sensor PIR selama 60 detik pertama menggunakan \texttt{time.monotonic()} sebelum mengizinkan pembacaan deteksi gerakan.
    \item \textbf{Akuisisi Sensor Tangguh (\texttt{baca\_semua\_sensor}):} Membaca nilai suhu dan kelembaban di dalam blok 
    \texttt{try-except RuntimeError} untuk mencegah program berhenti paksa ketika terjadi kegagalan \textit{checksum} pada DHT11, 
    serta membaca status gerakan PIR jika masa \textit{warm-up} telah selesai.
    \item \textbf{Mesin Evaluasi State (\texttt{evaluasi\_state}):} Menentukan transisi kondisi sistem berdasarkan hierarki 
    prioritas (Prioritas 1 untuk \texttt{CRITICAL}, Prioritas 2 untuk \texttt{WARNING}, dan kembali ke \texttt{NORMAL} sekaligus 
    me-reset status \texttt{is\_muted} secara otomatis jika kondisi kembali aman).
    \item \textbf{Pengendali Aktuator (\texttt{jalankan\_output}):} Mengatur pola nyala/kedip LED dan frekuensi bunyi 
    \textit{buzzer} secara \textit{non-blocking} di latar belakang (\texttt{background=True}) setiap kali terjadi perubahan 
    \textit{state}.
\end{enumerate}

\subsection{Cuplikan Kode}
Berikut adalah cuplikan fungsi evaluasi \textit{state} dan kontrol aktuator yang menjadi inti pengambilan keputusan pada sistem:

\begin{lstlisting}[language=Python, caption={Cuplikan Fungsi Evaluasi State dan Kontrol Output}]
def evaluasi_state(suhu, kelembaban, gerakan):
    global current_state, is_muted

    # Prioritas 1: kondisi kritis
    if suhu >= BATAS_SUHU_TINGGI or kelembaban >= BATAS_LEMBAB_KRITIS:
        state_baru = "CRITICAL"
    # Prioritas 2: kondisi peringatan
    elif kelembaban >= BATAS_LEMBAB_WASPADA or gerakan:
        state_baru = "WARNING"
    else:
        state_baru = "NORMAL"
        is_muted = False

    # Mute tidak mengubah kondisi sensor
    if state_baru in ("WARNING", "CRITICAL") and is_muted:
        current_state = "MUTED"
    else:
        current_state = state_baru

    return current_state
\end{lstlisting}

% 6. PENGUJIAN DAN HASIL
\section{Pengujian dan Hasil}
\subsection{Skenario Uji}
Pengujian dilakukan melalui enam skenario yang mencakup masa kalibrasi awal, kondisi lingkungan normal, tingkat peringatan 
bertingkat (waspada hingga kritis), deteksi gerakan, serta fungsi peredaman suara (\textit{mute}). Rekapitulasi skenario dan 
hasil uji ditunjukkan pada Tabel \ref{tab:hasil_uji}.

\begin{table}[h!]
\centering
\caption{Rekapitulasi Skenario dan Hasil Pengujian Sistem}
\vspace{0.2cm}
\resizebox{\textwidth}{!}{%
\begin{tabular}{@{}lllll@{}}
\toprule
\textbf{No.} & \textbf{Skenario Uji} & \textbf{Kondisi Input / Parameter Uji} & \textbf{Respons Output Sistem} & \textbf{Status} \\ \midrule
1 & Inisialisasi (\textit{Warm-up}) & Waktu jalan $<$ 60 dtk, Suhu $<$ 31$^\circ$C, RH $<$ 80\% & State \texttt{INITIALIZING}, LED Hijau ON, PIR diabaikan & Berhasil \\
2 & Kondisi Normal & Waktu $\ge$ 60 dtk, Suhu $<$ 31$^\circ$C, RH $<$ 75\%, PIR Aman & State \texttt{NORMAL}, LED Hijau ON, Buzzer OFF & Berhasil \\
3 & Kelembaban Waspada / Gerak & RH 75\%--79\% ATAU Sensor PIR mendeteksi gerakan & State \texttt{WARNING}, LED Kuning ON, Buzzer Beep Lambat & Berhasil \\
4 & Suhu / Kelembaban Kritis & Suhu $\ge$ 31$^\circ$C ATAU RH $\ge$ 80\% (termasuk saat \textit{warm-up}) & State \texttt{CRITICAL}, LED Merah Kedip Cepat, Buzzer Cepat & Berhasil \\
5 & Penekanan Tombol (Saat Alarm) & Push button ditekan saat state \texttt{WARNING} / \texttt{CRITICAL} & State \texttt{MUTED}, LED Kuning \& Merah ON, Buzzer OFF & Berhasil \\
6 & Penekanan Tombol (Saat Aman) & Push button ditekan saat state \texttt{NORMAL} / \texttt{INITIALIZING} & Pesan penolakan di terminal, state tidak berubah & Berhasil \\ \bottomrule
\end{tabular}%
}
\label{tab:hasil_uji}
\end{table}

\subsection{Analisis dan Hasil Pengujian}
Berdasarkan implementasi kode program dan skenario pengujian yang telah dijalankan, diperoleh analisis teknis sebagai berikut:

\begin{enumerate}
    \item \textbf{Kesesuaian Logika Prioritas dan Penanganan Masa Kalibrasi (\textit{Warm-up}):}
    Sensor PIR membutuhkan waktu stabilisasi termal sekitar 60 detik saat pertama kali dinyalakan agar tidak memicu alarm palsu 
    (\textit{false positive}). Kode program berhasil menangani hal ini melalui fungsi \texttt{periksa\_warmup\_pir()} yang menahan
     status sistem pada \texttt{INITIALIZING} (LED Hijau menyala) dan mengabaikan pembacaan PIR. Namun, program dilengkapi logika
      keselamatan cerdas: apabila selama masa \textit{warm-up} terjadi kondisi darurat (suhu $\ge$ 31.0$^\circ$C atau kelembaban 
      $\ge$ 80.0\%), sistem langsung memotong proses \textit{warm-up} dan masuk ke state \texttt{CRITICAL}.
    
    \item \textbf{Mekanisme Interupsi Tombol Mute dan Auto-Reset:}
    Penggunaan \textit{event-driven callback} (\texttt{tombol\_mute.when\_pressed}) dengan parameter \texttt{bounce\_time=0.1} 
    detik membuat respons tombol sangat instan tanpa harus menunggu jeda \texttt{time.sleep(2.0)} pada siklus utama selesai. 
    Ketika alarm diheningkan (\texttt{MUTED}), buzzer berhenti berbunyi tetapi LED Kuning dan Merah dinyalakan secara bersamaan 
    sebagai penanda visual bahwa bahaya belum berakhir. Saat kondisi ruangan kembali ke \texttt{NORMAL}, variabel 
    \texttt{is\_muted} otomatis diatur ulang menjadi \texttt{False}, sehingga alarm siap berbunyi kembali jika terjadi 
    bahaya berikutnya.

    \item \textbf{Kendala yang Muncul dan Solusi Perangkat Lunak:}
    \begin{itemize}
        \item \textbf{Kegagalan Baca Sensor DHT11 (\textit{RuntimeError}):} Sensor DHT11 menggunakan protokol komunikasi satu 
        jalur yang sangat sensitif terhadap waktu (\textit{timing}), sehingga kerap memunculkan eksepsi \texttt{RuntimeError} 
        (seperti kegagalan \textit{checksum} atau \textit{buffer} belum siap). Kendala ini berhasil diatasi dengan membungkus 
        pembacaan sensor di dalam blok \texttt{try-except RuntimeError} serta menyimpan nilai pembacaan terakhir yang valid ke 
        dalam variabel \texttt{suhu\_valid} dan \texttt{lembab\_valid} (dengan nilai awal *default* 27.0$^\circ$C dan 60.0\%).
        \item \textbf{Transisi Output Tanpa Kedipan Berulang (\textit{State Change Detection}):} Jika fungsi pengendali LED dan 
        buzzer dipanggil ulang pada setiap iterasi \textit{loop} (setiap 2 detik), pola kedipan (\texttt{blink}/\texttt{beep}) 
        akan terus ter-reset dari awal. Penambahan variabel pelacak \texttt{status\_output\_terakhir} memastikan fungsi 
        \texttt{jalankan\_output()} hanya dieksekusi satu kali saat terjadi perpindahan \textit{state}.
    \end{itemize}
\end{enumerate}

% 7. PEMBAHASAN
\section{Pembahasan}
\subsection{Evaluasi Performa Sistem}
Secara arsitektural, sistem bekerja dengan sangat efisien dan responsif. Pemanfaatan fitur \texttt{background=True} pada 
metode \texttt{blink()} dan \texttt{beep()} di pustaka \texttt{gpiozero} memungkinkan pembangkitan sinyal PWM/kedipan LED 
(siklus 0.2 detik) dan ketukan buzzer (siklus 0.1 detik pada kondisi kritis) berjalan pada \textit{thread} terpisah. Hal ini 
membuat \textit{main loop} tetap dapat berfokus pada pembacaan sensor setiap 2.0 detik tanpa memblokir frekuensi kedipan indikator.
 Selain itu, penggunaan \texttt{time.monotonic()} menjamin akurasi perhitungan waktu \textit{warm-up} karena nilainya tidak 
 terpengaruh apabila sistem operasi melakukan sinkronisasi waktu jaringan (NTP). Pada saat program dihentikan 
 (\texttt{KeyboardInterrupt}), blok \texttt{finally} juga memastikan seluruh objek perangkat keras ditutup (\texttt{close()} 
 dan \texttt{dht\_device.exit()}) sehingga tidak menyisakan pin GPIO yang menggantung.

\subsection{Kelebihan dan Kekurangan}
\begin{itemize}
    \item \textbf{Kelebihan:}
    \begin{enumerate}
        \item \textbf{Toleransi Kesalahan Tinggi (\textit{Fault-Tolerant}):} 
        Program tidak akan berhenti mendadak (\textit{crash}) saat sensor DHT11 gagal merespons berkat mekanisme penyimpanan 
        nilai valid terakhir (\textit{last known good value}).
        \item \textbf{Indikasi Bertingkat yang Informatif:} 
        Membedakan secara tegas antara peringatan dini (\texttt{WARNING} dengan LED kuning dan beep lambat) serta kondisi bahaya 
        (\texttt{CRITICAL} dengan LED merah berkedip cepat dan beep cepat).
        \item \textbf{Interaksi Tombol Non-Blocking:} 
        Penekanan tombol (\textit{mute}) diproses secara asinkron menggunakan \textit{event listener} dan dilengkapi proteksi 
        \textit{software debouncing} 100 ms.
    \end{enumerate}
    \item \textbf{Kekurangan:}
    \begin{enumerate}
        \item \textbf{Keterbatasan Akurasi DHT11:} Sensor DHT11 memiliki resolusi pembacaan yang terbatas dan rentang 
        \textit{sampling rate} maksimal 1 Hz (1 kali per detik), sehingga kurang ideal untuk mendeteksi lonjakan suhu mikro 
        secara instan.
        \item \textbf{Tidak Adanya Histeresis Ambang Batas:} Apabila suhu ruangan berfluktuasi tipis di sekitar angka ambang batas 
        (misalnya naik-turun antara 30.9$^\circ$C dan 31.0$^\circ$C), sistem dapat berpindah-pindah \textit{state} antara 
        \texttt{NORMAL} dan \texttt{CRITICAL} secara berulang (\textit{chattering}).
        \item \textbf{Penyimpanan Data Masih Volatil:} Log pembacaan sensor hanya ditampilkan ke layar terminal 
        (\textit{standard output}) dan belum disimpan ke dalam berkas log (\textit{file CSV}) maupun basis data.
    \end{enumerate}
\end{itemize}

\subsection{Potensi Pengembangan}
Untuk meningkatkan keandalan sistem di lingkungan ruang server atau laboratorium yang sesungguhnya, beberapa pengembangan yang dapat dilakukan meliputi:
\begin{enumerate}
    \item \textbf{Penambahan Logika Histeresis (\textit{Deadband}):} Menerapkan batas bawah pemulihan yang berbeda dari batas atas pemicu (contoh: masuk state \texttt{CRITICAL} saat suhu $\ge$ 31.0$^\circ$C, namun baru kembali ke \texttt{NORMAL} ketika suhu sudah turun hingga $\le$ 29.5$^\circ$C) untuk mencegah osilasi \textit{state}.
    \item \textbf{Pencatatan Data (\textit{Data Logging}) \& Integrasi IoT:} Menambahkan modul pencatatan otomatis ke basis data lokal (SQLite/InfluxDB) serta mengirimkan peringatan darurat secara \textit{real-time} melalui bot Telegram atau protokol MQTT.
    \item \textbf{Otomasi Aktuator Pendingin:} Menambahkan modul \textit{relay} untuk menyalakan kipas \textit{exhaust} atau unit AC cadangan secara otomatis ketika sistem memasuki kondisi \texttt{CRITICAL}.
\end{enumerate}

% 8. KESIMPULAN
\section{Kesimpulan}
Berdasarkan perancangan, implementasi, dan pengujian yang telah dilakukan, dapat disimpulkan bahwa:
\begin{enumerate}
    \item Sistem monitoring ruangan sederhana berbasis Raspberry Pi berhasil dirancang dan diimplementasikan dengan mengintegrasikan sensor suhu dan kelembaban (DHT), \textit{push button}, serta indikator output LED/Buzzer melalui pin GPIO.
    \item Perangkat lunak yang dibangun menggunakan bahasa Python mampu mengakuisisi data sensor secara berkala dan mengeksekusi transisi \textit{state machine} secara akurat sesuai ambang batas yang ditentukan.
    \item Berdasarkan empat skenario pengujian (kondisi normal, suhu tinggi, kelembaban tinggi, dan penekanan tombol), sistem mampu merespons perubahan kondisi lingkungan dan interaksi pengguna sesuai dengan tujuan proyek.
\end{enumerate}

\newpage
% 9. LAMPIRAN
\section*{Lampiran}
\addcontentsline{toc}{section}{Lampiran}

\subsection*{A. Kode Lengkap Program (\texttt{main.py})}
\begin{lstlisting}[language=Python, caption={Kode Lengkap Sistem Monitoring Ruangan Sederhana}]
import time
import board
import adafruit_dht


from gpiozero import (
    TrafficLights,
    Buzzer,
    Button,
    MotionSensor,
)


# ============================================================
# 1. INISIALISASI HARDWARE
# ============================================================


lampu = TrafficLights(
    red=17,
    yellow=27,
    green=22,
)


buzzer = Buzzer(25)


tombol_mute = Button(
    5,
    pull_up=True,
    bounce_time=0.1,
)


sensor_pir = MotionSensor(23)


dht_device = adafruit_dht.DHT11(
    board.D4,
    use_pulseio=False,
)


# ============================================================
# 2. KONFIGURASI SISTEM
# ============================================================


WAKTU_WARMUP_PIR = 60  # detik


BATAS_SUHU_TINGGI = 31.0
BATAS_LEMBAB_WASPADA = 75.0
BATAS_LEMBAB_KRITIS = 80.0


current_state = "INITIALIZING"
is_muted = False
pir_siap = False


# ============================================================
# 3. TOMBOL MUTE
# ============================================================


def tekan_tombol_handler():
    global is_muted


    if current_state in ("WARNING", "CRITICAL", "MUTED"):
        is_muted = not is_muted


        status_teks = (
            "DIHENINGKAN (MUTED)"
            if is_muted
            else "AKTIF KEMBALI"
        )


        print(
            f"\n>>> [BUTTON DITEKAN] "
            f"Alarm Buzzer: {status_teks} <<<"
        )
    else:
        print(
            "\n>>> [BUTTON DITEKAN] "
            "Sistem belum dalam kondisi alarm. <<<"
        )




tombol_mute.when_pressed = tekan_tombol_handler


# ============================================================
# 4. WARM-UP PIR
# ============================================================


def periksa_warmup_pir(waktu_mulai):
    global pir_siap


    if pir_siap:
        return


    waktu_berjalan = time.monotonic() - waktu_mulai
    sisa = max(
        0,
        int(WAKTU_WARMUP_PIR - waktu_berjalan),
    )


    if waktu_berjalan >= WAKTU_WARMUP_PIR:
        pir_siap = True
        print(
            "\n[PIR SIAP] Sensor gerakan siap digunakan.\n"
        )
    else:
        print(
            f"\r[PIR WARM-UP] {sisa:02d} detik tersisa",
            end="",
            flush=True,
        )




# ============================================================
# 5. PEMBACAAN SENSOR
# ============================================================


def baca_semua_sensor():
    suhu = None
    kelembaban = None


    try:
        suhu = dht_device.temperature
        kelembaban = dht_device.humidity
    except RuntimeError:
        # Jika pembacaan gagal, coba lagi pada iterasi berikutnya.
        pass


    ada_gerakan = (
        sensor_pir.motion_detected
        if pir_siap
        else False
    )


    return suhu, kelembaban, ada_gerakan




# ============================================================
# 6. EVALUASI STATUS
# ============================================================


def evaluasi_state(suhu, kelembaban, gerakan):
    global current_state, is_muted


    # Prioritas 1: kondisi kritis
    if (
        suhu >= BATAS_SUHU_TINGGI
        or kelembaban >= BATAS_LEMBAB_KRITIS
    ):
        state_baru = "CRITICAL"


    # Prioritas 2: kondisi peringatan
    elif (
        kelembaban >= BATAS_LEMBAB_WASPADA
        or gerakan
    ):
        state_baru = "WARNING"


    else:
        state_baru = "NORMAL"
        is_muted = False


    # Mute tidak mengubah kondisi sensor.
    if state_baru in ("WARNING", "CRITICAL") and is_muted:
        current_state = "MUTED"
    else:
        current_state = state_baru


    return current_state




# ============================================================
# 7. KONTROL LAMPU DAN BUZZER
# ============================================================


def jalankan_output(state):
    # Hentikan pola output sebelumnya.
    buzzer.off()
    lampu.off()


    if state == "INITIALIZING":
        lampu.green.on()


    elif state == "NORMAL":
        lampu.green.on()


    elif state == "WARNING":
        lampu.yellow.on()


        buzzer.beep(
            on_time=0.2,
            off_time=0.8,
            background=True,
        )


    elif state == "CRITICAL":
        lampu.red.blink(
            on_time=0.2,
            off_time=0.2,
            background=True,
        )


        buzzer.beep(
            on_time=0.1,
            off_time=0.1,
            background=True,
        )


    elif state == "MUTED":
        lampu.yellow.on()
        lampu.red.on()
        buzzer.off()




# ============================================================
# 8. PROGRAM UTAMA
# ============================================================


def main():
    global current_state, pir_siap


    print("=== SISTEM MONITORING RUANG LAB/SERVER AKTIF ===")


    # Nilai awal sementara jika DHT11 belum berhasil dibaca.
    suhu_valid = 27.0
    lembab_valid = 60.0


    waktu_mulai = time.monotonic()
    status_output_terakhir = "INITIALIZING"


    # LED hijau sebagai indikator inisialisasi.
    jalankan_output("INITIALIZING")


    try:
        while True:
            periksa_warmup_pir(waktu_mulai)


            # DHT11 tetap dibaca selama warm-up PIR.
            suhu, kelembaban, gerakan = baca_semua_sensor()


            if suhu is not None and kelembaban is not None:
                suhu_valid = suhu
                lembab_valid = kelembaban


            # Selama warm-up, hijau tetap menyala kecuali
            # kondisi kritis suhu atau kelembapan terjadi.
            if not pir_siap:
                if (
                    suhu_valid >= BATAS_SUHU_TINGGI
                    or lembab_valid >= BATAS_LEMBAB_KRITIS
                ):
                    state_aktif = evaluasi_state(
                        suhu_valid,
                        lembab_valid,
                        False,
                    )
                else:
                    state_aktif = "INITIALIZING"
                    current_state = state_aktif
            else:
                state_aktif = evaluasi_state(
                    suhu_valid,
                    lembab_valid,
                    gerakan,
                )


            # Ubah output hanya saat status berubah.
            if state_aktif != status_output_terakhir:
                jalankan_output(state_aktif)
                status_output_terakhir = state_aktif


            if not pir_siap:
                status_pir = "WARM-UP"
            elif gerakan:
                status_pir = "GERAK"
            else:
                status_pir = "AMAN"


            print(
                f"[STATE: {state_aktif:12s}] "
                f"Suhu: {suhu_valid:.1f}C | "
                f"Lembab: {lembab_valid}% | "
                f"PIR: {status_pir:7s} | "
                f"Mute: {is_muted}"
            )


            time.sleep(2.0)


    except KeyboardInterrupt:
        print(
            "\nMenghentikan program dan membersihkan pin GPIO..."
        )


    finally:
        buzzer.off()
        lampu.off()


        tombol_mute.close()
        sensor_pir.close()


        dht_device.exit()


        buzzer.close()
        lampu.close()


        print("Semua output dimatikan dan GPIO dibersihkan.")




if __name__ == "__main__":
    main()

\end{lstlisting}

\subsection*{B. Dokumentasi Foto}
Berikut adalah dokumentasi rangkaian perangkat keras dan pengujian sistem:

\vspace{0.3cm}
\noindent\framebox[\textwidth]{%
\parbox{0.95\textwidth}{%
\centering
\vspace{1.2cm}
\textbf{[Tempat Foto 1: Rangkaian Keseluruhan Sistem Monitoring Ruangan]}\\
\small (Deskripsi: Rangkaian Raspberry Pi, Sensor DHT, Push Button, dan LED pada Breadboard)
\vspace{1.2cm}
}}

\vspace{0.4cm}
\noindent\framebox[\textwidth]{%
\parbox{0.95\textwidth}{%
\centering
\vspace{1.2cm}
\textbf{[Tempat Foto 2: Pengujian Saat Kondisi Normal dan Kondisi Peringatan]}\\
\small (Deskripsi: Tampilan output terminal beserta indikator LED menyala)
\vspace{1.2cm}
}}

\subsection*{C. Link Video Demo Pembuatan dan Pengujian}
Video demonstrasi perakitan, penjelasan alur program, serta pengujian langsung sistem dapat diakses melalui tautan berikut:
\begin{itemize}
    \item \textbf{Tautan (URL):} \url{https://vt.tiktok.com/ZSgaYhkMV/}
\end{itemize}

\end{document}